\section{Attitude Control and Quaternion Optimization}
\label{sec:app-attitude-control-and-quaternion-optimization}

\begin{remark}
	Many robotic optimal-control problems live on nonlinear configuration spaces rather than Euclidean spaces. Attitude is the simplest and most important example: the physically meaningful state is a rotation, not an arbitrary vector in $\mathbb{R}^{3\times 3}$ or $\mathbb{R}^{4}$. This appendix summarizes the geometric constructions most often needed when trajectory optimization, estimation, and LQR are carried out with attitude states.
\end{remark}

\subsection{Rotation Matrices and Unit Quaternions}
\label{sec:app-rotation-matrices-and-unit-quaternions}

\begin{definition}[Special orthogonal group]\label{def:app-special-orthogonal-group}
	The set of three-dimensional rotation matrices is
	\[
		SO(3)
		=
		\left\{
			R \in \mathbb{R}^{3\times 3}
			:
			R^\top R = I,
			\ \det R = 1
		\right\}.
	\]
\end{definition}

\begin{definition}[Hat map]\label{def:app-hat-map}
	For $\omega=(\omega_1,\omega_2,\omega_3)^\top \in \mathbb{R}^{3}$, define
	\[
		\widehat{\omega}
		=
		\begin{bmatrix}
			0          & -\omega_3 & \omega_2 \\
			\omega_3   & 0         & -\omega_1 \\
			-\omega_2  & \omega_1  & 0
		\end{bmatrix}.
	\]
	Then
	\[
		\widehat{\omega}v=\omega \times v
	\]
	for every $v \in \mathbb{R}^{3}$.
\end{definition}

\begin{definition}[Unit quaternion]\label{def:app-unit-quaternion}
	A unit quaternion is a vector
	\[
		q=
		\begin{bmatrix}
			q_0 \\
			q_v
		\end{bmatrix}
		\in \mathbb{R}^{4},
		\qquad
		q_0 \in \mathbb{R},
		\quad
		q_v \in \mathbb{R}^{3},
	\]
	satisfying
	\[
		q^\top q = q_0^2+\|q_v\|^2 = 1.
	\]
	The set of unit quaternions is the three-sphere
	\[
		S^3=\{q \in \mathbb{R}^{4}:q^\top q=1\}.
	\]
\end{definition}

\begin{remark}[Double cover]
	Every rotation in $SO(3)$ is represented by two antipodal unit quaternions:
	\[
		q
		\qquad \text{and} \qquad
		-q.
	\]
	They encode the same physical attitude. This is not a singularity, but it does matter numerically because optimization and filtering routines should maintain a consistent sign convention along a trajectory.
\end{remark}

\begin{definition}[Quaternion-to-rotation map]\label{def:app-quaternion-to-rotation-map}
	The rotation matrix corresponding to a unit quaternion
	\[
		q=
		\begin{bmatrix}
			q_0 \\
			q_v
		\end{bmatrix}
	\]
	is
	\[
		R(q)
		=
		(q_0^2-q_v^\top q_v)I
		+
		2q_v q_v^\top
		+
		2q_0 \widehat{q_v}.
	\]
\end{definition}

\begin{definition}[Quaternion kinematics]\label{def:app-quaternion-kinematics}
	If $\omega \in \mathbb{R}^{3}$ is the body angular velocity, then the quaternion kinematics are
	\[
		\dot q = \frac{1}{2}G(q)\omega,
	\]
	where
	\[
		G(q)
		=
		\begin{bmatrix}
			-q_v^\top \\
			q_0 I + \widehat{q_v}
		\end{bmatrix}
		\in \mathbb{R}^{4\times 3}.
	\]
\end{definition}

\begin{proposition}[Tangent-space property of the quaternion kinematics matrix]\label{prop:app-tangent-space-property-of-the-quaternion-kinematics-matrix}
	For every unit quaternion $q \in S^3$,
	\[
		q^\top G(q)=0
	\]
	and
	\[
		G(q)^\top G(q)=I_3.
	\]
	Therefore the columns of $G(q)$ form an orthonormal basis for the tangent space
	\[
		T_q S^3 = \{\delta q \in \mathbb{R}^{4}:q^\top \delta q=0\}.
	\]
\end{proposition}
\begin{proof}
	Direct multiplication gives
	\[
		q^\top G(q)
		=
		\begin{bmatrix}
			q_0 & q_v^\top
		\end{bmatrix}
		\begin{bmatrix}
			-q_v^\top \\
			q_0 I + \widehat{q_v}
		\end{bmatrix}
		=
		-q_0 q_v^\top + q_0 q_v^\top + q_v^\top \widehat{q_v}
		=
		0,
	\]
	because
	\[
		q_v^\top \widehat{q_v}=0.
	\]
	A second direct multiplication yields
	\[
		G(q)^\top G(q)
		=
		q_v q_v^\top + (q_0 I-\widehat{q_v})(q_0 I+\widehat{q_v})
		=
		(q_0^2+\|q_v\|^2)I
		=
		I_3,
	\]
	using
	\[
		\widehat{q_v}^{\,2}=q_v q_v^\top-\|q_v\|^2 I
	\]
	and the unit-norm condition on $q$.
\end{proof}

\subsection{Optimization Updates in the Quaternion Tangent Space}
\label{sec:app-optimization-updates-in-the-quaternion-tangent-space}

\begin{remark}
	An unconstrained Newton or Gauss--Newton method wants to add a correction vector to the current iterate. For quaternions this is invalid: an arbitrary
	\[
		q+\delta q
	\]
	will not remain on $S^3$. The correction must instead be computed in a three-dimensional local coordinate chart and then mapped back to the manifold.
\end{remark}

\begin{definition}[Small-angle quaternion update]\label{def:app-small-angle-quaternion-update}
	Let $\phi \in \mathbb{R}^{3}$ be a small attitude correction. A corresponding first-order unit-quaternion update is
	\[
		\delta q(\phi)
		\approx
		\begin{bmatrix}
			1 \\
			\frac{1}{2}\phi
		\end{bmatrix},
	\]
	and the updated quaternion is obtained multiplicatively,
	\[
		q^{+} = \delta q(\phi)\otimes q,
	\]
	followed by normalization if needed.
\end{definition}

\begin{remark}[Why multiplicative updates are natural]
	Additive updates live in the ambient space $\mathbb{R}^{4}$. Multiplicative updates live on the group structure of rotations. The local coordinate $\phi \in \mathbb{R}^{3}$ is unconstrained, while the updated quaternion remains on the unit sphere after renormalization.
\end{remark}

\begin{definition}[Full-state projection matrix for error coordinates]\label{def:app-full-state-projection-matrix-for-error-coordinates}
	Suppose the state contains position, quaternion attitude, linear velocity, and angular velocity:
	\[
		x=
		\begin{bmatrix}
			r \\
			q \\
			v \\
			\omega
		\end{bmatrix}
		\in \mathbb{R}^{13}.
	\]
	The tangent-space error coordinate
	\[
		\eta=
		\begin{bmatrix}
			\delta r \\
			\delta \phi \\
			\delta v \\
			\delta \omega
		\end{bmatrix}
		\in \mathbb{R}^{12}
	\]
	is lifted to an ambient perturbation by
	\[
		\delta x = E(x)\eta,
	\]
	with
	\[
		E(x)
		=
		\operatorname{diag}\!\bigl(I_3,G(q),I_3,I_3\bigr)
		\in \mathbb{R}^{13\times 12}.
	\]
\end{definition}

\begin{proposition}[Projected linearization in quaternion error coordinates]\label{prop:app-projected-linearization-in-quaternion-error-coordinates}
	Let
	\[
		x_{k+1}=f_k(x_k,u_k)
	\]
	and linearize in the ambient coordinates:
	\[
		\delta x_{k+1}=A_k \delta x_k + B_k \delta u_k.
	\]
	If $\eta_k$ and $\eta_{k+1}$ are the tangent-space errors defined by
	\[
		\delta x_k = E(\bar x_k)\eta_k,
		\qquad
		\eta_{k+1}=E(\bar x_{k+1})^\top \delta x_{k+1},
	\]
	then the projected error dynamics are
	\[
		\eta_{k+1}
		=
		\underbrace{E(\bar x_{k+1})^\top A_k E(\bar x_k)}_{A_k^\eta}\eta_k
		+
		\underbrace{E(\bar x_{k+1})^\top B_k}_{B_k^\eta}\delta u_k.
	\]
\end{proposition}
\begin{proof}
	Substitute
	\[
		\delta x_k=E(\bar x_k)\eta_k
	\]
	into the ambient linearization:
	\[
		\delta x_{k+1}=A_k E(\bar x_k)\eta_k+B_k \delta u_k.
	\]
	Project the result back onto the tangent space at the next nominal state by multiplying with
	\[
		E(\bar x_{k+1})^\top.
	\]
	This gives the stated reduced linear dynamics directly.
\end{proof}

\subsection{Wahba's Problem and Gauss--Newton on $S^3$}
\label{sec:app-wahba-s-problem-and-gauss-newton-on-s3}

\begin{definition}[Wahba's problem]\label{def:app-wahba-s-problem}
	Given known inertial-frame vectors
	\[
		v_1^N,\dots,v_m^N \in \mathbb{R}^{3}
	\]
	and their corresponding body-frame measurements
	\[
		v_1^B,\dots,v_m^B \in \mathbb{R}^{3},
	\]
	find the attitude $q$ minimizing
	\[
		J(q)
		=
		\frac{1}{2}\sum_{i=1}^{m} w_i
		\left\|
			R(q)v_i^B-v_i^N
		\right\|^2,
	\]
	where $w_i>0$ are weights.
\end{definition}

\begin{remark}
	Wahba's problem is the canonical example of optimization on a rotation manifold. It is also a clean model for attitude estimation from vector observations, such as star trackers, accelerometers, or magnetometers.
\end{remark}

\begin{algorithm}[H]
	\DontPrintSemicolon
	\caption{Gauss--Newton for Wahba's Problem in Quaternion Coordinates}
\label{alg:app-gauss-newton-for-wahba-s-problem-in-quaternion-coordinates}
	\KwIn{Initial unit quaternion $q_0$, vector correspondences $(v_i^B,v_i^N)$, weights $w_i$}
	\For{$r=0,1,2,\dots$ until convergence}{
		Compute residuals
		\[
			r_i(q_r)=R(q_r)v_i^B-v_i^N
		\]
		and stack them into $r(q_r)$\;
		Linearize the residual with respect to a tangent-space correction $\phi \in \mathbb{R}^{3}$:
		\[
			r(q_r \oplus \phi)\approx r(q_r)+J_\phi(q_r)\phi
		\]\;
		Solve the normal equations
		\[
			J_\phi(q_r)^\top J_\phi(q_r)\,\phi
			=
			- J_\phi(q_r)^\top r(q_r)
		\]\;
		Form the multiplicative update $q_{r+1}=\delta q(\phi)\otimes q_r$\;
		Renormalize $q_{r+1}$ to unit length\;
	}
	\KwOut{Estimated attitude quaternion}
\end{algorithm}

\begin{remark}[Role of the attitude Jacobian]
	The Jacobian
	\[
		J_\phi(q)
	\]
	is the chain-rule Jacobian of the residual with respect to the three-dimensional correction variable $\phi$, not with respect to the ambient four-dimensional quaternion coordinate. This is exactly the same philosophy as \autoref{prop:app-projected-linearization-in-quaternion-error-coordinates}: compute in the tangent space, then lift back to the manifold.
\end{remark}

\subsection{Quaternion Error-State LQR and Attitude Control}
\label{sec:app-quaternion-error-state-lqr-and-attitude-control}

\begin{definition}[Quaternion error state]\label{def:app-quaternion-error-state}
	Given a nominal quaternion $\bar q$ and actual quaternion $q$, define the error quaternion by
	\[
		\delta q = \bar q^\ast \otimes q,
	\]
	where $\bar q^\ast$ is the quaternion conjugate. For small errors,
	\[
		\delta q
		\approx
		\begin{bmatrix}
			1 \\
			\frac{1}{2}\delta \phi
		\end{bmatrix},
	\]
	so the three-vector $\delta \phi$ serves as the local attitude error coordinate.
\end{definition}

\begin{remark}[Why naive LQR on the ambient quaternion is wrong]
	A state containing a quaternion lives on a manifold of codimension one. A naive linearization in the ambient coordinate $\delta q \in \mathbb{R}^{4}$ creates a fictitious radial direction that changes $\|q\|$, but that direction is not physically actuated by the true dynamics. An LQR design built on the ambient linearization therefore solves the wrong local problem.
\end{remark}

\begin{algorithm}[H]
	\DontPrintSemicolon
	\caption{Quaternion Error-State LQR Design}
\label{alg:app-quaternion-error-state-lqr-design}
	\KwIn{Nominal attitude trajectory $\{\bar x_k,\bar u_k\}$ or equilibrium, ambient linearization matrices $(A_k,B_k)$, quadratic weights on tangent errors}
	Construct the projection matrices $E(\bar x_k)$ from \autoref{def:app-full-state-projection-matrix-for-error-coordinates}\;
	Project the ambient linearization onto tangent coordinates using \autoref{prop:app-projected-linearization-in-quaternion-error-coordinates}\;
	Choose a quadratic cost in the reduced error state
	\[
		\eta_k=
		\begin{bmatrix}
			\delta r_k \\
			\delta \phi_k \\
			\delta v_k \\
			\delta \omega_k
		\end{bmatrix}
	\]\;
	Solve the resulting reduced LQR problem using \autoref{alg:l6-finite-horizon-discrete-time-lqr} or the algebraic Riccati equation from \autoref{thm:l6-discrete-time-algebraic-riccati-equation-and-optimal-gain}\;
	Implement the feedback in quaternion error coordinates and renormalize the propagated quaternion state\;
	\KwOut{Locally optimal feedback law respecting the attitude manifold}
\end{algorithm}

\begin{remark}[Quadrotor attitude stabilization]
	For a rigid body or quadrotor, the practical pattern is standard:
	\begin{itemize}
		\item propagate the full nonlinear quaternion dynamics,
		\item compute the attitude error through the quaternion product,
		\item feed the reduced error vector into an LQR or MPC law designed in tangent coordinates,
		\item renormalize the quaternion after numerical integration.
	\end{itemize}
	This is the manifold-aware version of linear quadratic feedback. The control law is linear in the local error coordinates, not in the raw quaternion itself.
\end{remark}
